\documentclass[../../project_master.tex]{subfiles}
\begin{document}
\section{Interview Talk Outline}
The talk should be structured as a clear chain rather than as disconnected technical facts.

\subsection*{Opening Message}
The central message is that radiation damage in silicon can be understood as a multiscale physics problem: particle interactions deposit energy and create atomic disorder, that disorder introduces electrically active defects, and those defects degrade measurable device performance.

\subsection*{Suggested Narrative}
\begin{enumerate}[leftmargin=1.5em]
    \item Explain why silicon radiation damage matters in space, defense, and high-reliability electronics.
    \item Distinguish ionization from displacement damage.
    \item Show how Geant4 models the interaction stage inside silicon.
    \item Explain how recoil events and displacement energy transfer motivate an effective defect-density model.
    \item Show how defects shorten carrier lifetime and reduce diffusion length.
    \item Show how device observables such as leakage current worsen with exposure.
    \item End by emphasizing that the framework connects microscopic physics to engineering decisions.
\end{enumerate}

\subsection*{Tone to Use}
The strongest tone is technically honest and modest: Geant4 is presented as the radiation-transport engine, while MATLAB provides a reduced but physically interpretable solid-state model. This avoids overclaiming while still demonstrating strong interdisciplinary understanding.
\end{document}
